\section{Webinterface}

\subsection{Auswahl der Komponenten}

\subsubsection{Programmiersprachen}
Es stand am Anfang die Wahl zwischen PyScript und JavaScript. PyScript ist ein Tool, um Python 
in eine Web-Anwendung einzubinden, jedoch steht es noch in den Startlöchern und ist damit noch 
zu unprofessionell und unzuverlässig. Wäre PyScript weiterentwickelt, wäre es eine interessante 
Alternative zu JavaScript gewesen, da Python schon in den letzten Jahren gelernt wurde und 
JavaScript nicht. 

\subsubsection{WebTools für den Server}
Hier steht die Auswahl zwischen den Webservern Nginx und Apache im Vordergrund. Nginx ist die 
modernere und leistungsfähigere Variante von den beiden. Auch wenn momentan noch nicht viele 
Spieler gleichzeitig spielen, kann das in Zukunft vorkommen. Dazu ist ein leistungsfähiger und 
schneller Webserver wie Nginx essenziell. Hierfür wurde der Anwendungsserver uWSGi verwendet, 
da er sehr gut mit Nginx und Flask zusammenarbeitet \cite{FlaskDeploying}. 

Das Web-Framework Flask wurde gewählt, da es sehr minimalistisch startet und man selbst wählen 
kann, welche Bibliotheken man nutzt. Das führt dazu, dass weniger Arbeitsspeicher verwendet wird, 
was auf einem Debian-Server wichtig ist. Weiters braucht man durch die mit HTML und JavaScript 
geschriebene Website nur einen Socket, der PGN-Dateien empfängt. Flask ist genau darauf 
spezialisiert. Außerdem ist Flask sehr ähnlich zu Python, weshalb es leicht verständlich ist, 
es ist egal welche Datenbank man wählt und es arbeitet sehr nahe mit uWSGI und Nginx zusammen. 

\subsubsection{Server}
Es wurde ein Debian Server von der Schule verwendet. Er wird mit ssh konfiguriert. 

\subsubsection{Datenbank}
Die Wahl zwischen MariaDB und MySQL war nicht schwer. Das liegt daran, dass MariaDB eine 
modifizierte und damit modernere Version von MySQL ist. 

\subsection{Frontend}
Es soll eine webbasierte Benutzeroberfläche mit HTML, CSS und JavaScript erstellt werden, 
die der Benutzerin oder dem Benutzer den Zugriff auf seine Partien ermöglicht.

\subsection{Backend}
Ein Python-basiertes Flask-Framework, das als zentrale Schnittstelle zwischen der Datenbank und 
dem Frontend fungiert. Dieses verarbeitet eingehende Daten vom Frontend sowie Partiedaten vom 
physischen Schachbrett. Um parallele Prozesse (Web-Anfragen und Socket-Datenströme) zu bewältigen, 
wird die Multithreading-Funktion von Python genutzt.

\subsection{Datenbank}
Eine MariaDB-Datenbank, die alle benutzer- und spielrelevanten Daten (z.B. PGN-Notationen) 
speichert und über den SQL-Standard zur Verfügung stellt. 
